% B19/A094 independent reconstruction of B01 two-surgery refund.
\section*{Independent review: the two-surgery difference refund}

All logarithms are base two.  Put $\delta=2^{-k}$.  For distinct
$x,y\in\{0,1\}^k$, let $z=x\oplus y$, $R=\operatorname{supp}(z)$,
$r=|R|$, $\tau=g(x)g(y)$, and $\sigma=\tau(-1)^r$.  The two point-flip
Walsh perturbations are
\[
 e_x(T)=-2g(x)\delta\chi_T(x),\qquad
 e_y(T)=-2g(y)\delta\chi_T(y).
\]
For $D_e(S)=\sum_{T\supseteq S}e(T)^2$, write $T=S\cup W$.  Then
\[
 \sum_{T\supseteq S}\chi_T(z)
 =\chi_S(z)\prod_{j\notin S}(1+(-1)^{z_j})
 =2^{k-|S|}(-1)^r\mathbf 1_{R\subseteq S}.
\]
Consequently
\[
 D_{e_x+e_y}(S)
 =8\delta2^{-|S|}(1+\sigma\mathbf 1_{R\subseteq S})
\]
and, because $D_{e_x}(S)=D_{e_y}(S)=4\delta2^{-|S|}$,
\[
 C_R(S)=8\sigma\delta2^{-|S|}\mathbf 1_{R\subseteq S}.
\]
Under the flat bent law $p(S)=\delta$,
\[
 \sum_Sp(S)\frac{|C_R(S)|}{2^{-|S|}}
 =8\delta\,2^{-r}.
\]

Now define
\[
 \mathcal U(e)=\sum_Te(T)^2\log\frac{D_e(T)}{e(T)^2}.
\]
The combined atom law is
\[
 q(T)=(e_x(T)+e_y(T))^2
 =8\delta^2(1+\tau\chi_T(z)).
\]
There are $2^{k-1}$ nonzero atoms, each equal to $16\delta^2$.  On a
nonzero atom,
\[
 \log\frac{D_{e_x+e_y}(T)}{q(T)}
 =k-|T|-1+\mathbf 1_{R\subseteq T}.
\]
The final indicator can occur only when $\sigma=1$.

If $r\ge2$, conditioning on the parity of $|T\cap R|$ preserves every
coordinate marginal at $1/2$.  Hence
\[
 \sum_{T:q(T)>0}(k-|T|-1)=2^{k-2}(k-2),
\]
and the number of nonzero atoms containing $R$ is
$\mathbf 1_{\sigma=1}2^{k-r}$.  Therefore
\[
 \mathcal U(e_x+e_y)
 =4\delta(k-2)+16\delta2^{-r}\mathbf 1_{\sigma=1}.
\]
Since $\mathcal U(e_x)=\mathcal U(e_y)=2\delta k$, the refund is
\[
 \rho_{r,\sigma}=8\delta-2^{4-r}\delta\mathbf 1_{\sigma=1}
 \qquad(r\ge2).
\]

If $r=1$, the parity condition fixes the unique bit in $R$.  The
$\mathbf 1_{R\subseteq T}$ correction cancels this fixed membership, so
\[
 \mathcal U(e_x+e_y)=4\delta(k-1),\qquad \rho_{1,\sigma}=4\delta.
\]
It follows in all cases that
\[
 4\delta\le\rho_{r,\sigma}\le8\delta,
 \qquad 8\delta2^{-r}\le\rho_{r,\sigma}.
\]

These equations accept the two algebraic v1 claims.  Bentness is only
needed to interpret $p(S)=\delta$ as the flat parent spectral law; the
character-sum and difference-vector identities themselves hold for either
choice of the two phase signs.  The principal-filter generator $R$ is an
algebraic owner, not yet a disjoint physical pivotal event.  A complete
selector payment must restore the selector probability multiplying the
difference packet and bound the coherent-parent and retained-child
residual at the same normalization.  No such step is proved here.
